% GENERATED FILE. Edit canonical lessons, then run npm run generate:books.
% canonical-source-hash: fnv1a64:6a234c1f67cbc25f
% canonical-lessons: GU-C187-vadhvu, GU-C187-ghatvu, GU-C187-vadhare, GU-C187-tako, GU-C187-chut

\chapter{Prices that rise and fall}
\label{ch:gu-prices-that-rise-and-fall}
\hlchaptermodality{187}

\begin{chapteropening}
\textbf{By the end of this chapter:} \emph{I can read whether prices go up or down and what is on offer.}

Complete the last lesson of chapter 187: I can read whether prices go up or down and what is on offer.
\end{chapteropening}

\section[\texorpdfstring{\gu{વધવું} (vadhvũ)}{વધવું (vadhvũ)}]{\gu{વધવું} (vadhvũ) — to go up, to grow}

\label{lesson:GU-C187-vadhvu}



\begin{quote}
Before the new one: say the Gujarati for banknote, then the Gujarati for receipt.
\end{quote}

\subsection*{You'll want to know: \gu{વધવું}}
\textbf{\gu{વધવું}} — \emph{vadhvũ} — \textquotedblleft{}to go up, to grow\textquotedblright{}.

\begin{quote}
\textbf{\gu{દૂધનો} \gu{ભાવ} \gu{વધે} \gu{છે}.} — \emph{dūdhno bhāv vadhe che} — \textquotedblleft{}The price of milk is going up.\textquotedblright{}
\end{quote}

\begin{cousinweb}
\textbf{\gu{વધવું}} (\emph{vadhvũ}) is a tadbhava. CDIAL lists it, glossed \textquotedblleft{}to advance\textquotedblright{}, under Sanskrit \emph{várdhatē}, \textquotedblleft{}grows, increases\textquotedblright{} (CDIAL 11376).
\end{cousinweb}

\begin{grammarlens}[title={What you've built}]
One more word for everyday life at A2.
\end{grammarlens}

\subsection*{Your turn}
\begin{itemize}
\raggedright
  \item \emph{Say it:} \emph{vadhvũ}
  \item \emph{Say it:} \emph{vadhvũ}, once more
  \item \emph{Say it:} \emph{dūdhno bhāv vadhe che} — The price of milk is going up.
  \item \emph{Recall:} say the Gujarati for banknote, then the Gujarati for receipt, then say \emph{vadhvũ} again
\end{itemize}

\begin{tcolorbox}[breakable,colback=teal!4,colframe=teal!35!black,title={Before you move on}]
What does \gu{વધવું} mean? (\textquotedblleft{}To go up, to grow\textquotedblright{}.) Say it once more.
\end{tcolorbox}
\section[\texorpdfstring{\gu{ઘટવું} (ghaṭvũ)}{ઘટવું (ghaṭvũ)}]{\gu{ઘટવું} (ghaṭvũ) — to go down, to decrease}

\label{lesson:GU-C187-ghatvu}



\begin{quote}
Before the new one: say the Gujarati for receipt, then the Gujarati for to go up.
\end{quote}

\subsection*{You'll want to know: \gu{ઘટવું}}
\textbf{\gu{ઘટવું}} — \emph{ghaṭvũ} — \textquotedblleft{}to go down, to decrease\textquotedblright{}.

\begin{quote}
\textbf{\gu{શિયાળામાં} \gu{શાકનો} \gu{ભાવ} \gu{ઘટે} \gu{છે}.} — \emph{śiyāḷāmā̃ śākno bhāv ghaṭe che} — \textquotedblleft{}In winter the price of vegetables goes down.\textquotedblright{}
\end{quote}

\begin{cousinweb}
\textbf{\gu{ઘટવું}} (\emph{ghaṭvũ}) is a tadbhava. It comes from \emph{*ghaṭṭati}, \textquotedblleft{}decreases, is wanting\textquotedblright{}, a reconstructed form, and CDIAL glosses the Gujarati verb \textquotedblleft{}to decrease\textquotedblright{} (CDIAL 4415).
\end{cousinweb}

\begin{grammarlens}[title={What you've built}]
One more word for everyday life at A2.
\end{grammarlens}

\subsection*{Your turn}
\begin{itemize}
\raggedright
  \item \emph{Say it:} \emph{ghaṭvũ}
  \item \emph{Say it:} \emph{ghaṭvũ}, once more
  \item \emph{Say it:} \emph{śiyāḷāmā̃ śākno bhāv ghaṭe che} — In winter the price of vegetables goes down.
  \item \emph{Recall:} say the Gujarati for receipt, then the Gujarati for to go up, then say \emph{ghaṭvũ} again
\end{itemize}

\begin{tcolorbox}[breakable,colback=teal!4,colframe=teal!35!black,title={Before you move on}]
What does \gu{ઘટવું} mean? (\textquotedblleft{}To go down, to decrease\textquotedblright{}.) Say it once more.
\end{tcolorbox}
\section[\texorpdfstring{\gu{વધારે} (vadhāre)}{વધારે (vadhāre)}]{\gu{વધારે} (vadhāre) — more, too much}

\label{lesson:GU-C187-vadhare}



\begin{quote}
Before the new one: say the Gujarati for to go up, then the Gujarati for to go down.
\end{quote}

\subsection*{You'll want to know: \gu{વધારે}}
\textbf{\gu{વધારે}} — \emph{vadhāre} — \textquotedblleft{}more, too much\textquotedblright{}.

\begin{quote}
\textbf{\gu{આ} \gu{દુકાનમાં} \gu{વધારે} \gu{ગ્રાહક} \gu{છે}.} — \emph{ā dukānmā̃ vadhāre grāhak che} — \textquotedblleft{}This shop has more customers.\textquotedblright{}
\end{quote}

\begin{cousinweb}
\textbf{\gu{વધારે}} (\emph{vadhāre}) is a tadbhava. It comes from Sanskrit \emph{vṛddhatara}, \textquotedblleft{}larger\textquotedblright{}, the comparative of \emph{vṛddha}, \textquotedblleft{}grown\textquotedblright{}, and CDIAL glosses the Gujarati word \textquotedblleft{}more, greater\textquotedblright{} (CDIAL 12074).
\end{cousinweb}

\begin{grammarlens}[title={What you've built}]
One more word for everyday life at A2.
\end{grammarlens}

\subsection*{Your turn}
\begin{itemize}
\raggedright
  \item \emph{Say it:} \emph{vadhāre}
  \item \emph{Say it:} \emph{vadhāre}, once more
  \item \emph{Say it:} \emph{ā dukānmā̃ vadhāre grāhak che} — This shop has more customers.
  \item \emph{Recall:} say the Gujarati for to go up, then the Gujarati for to go down, then say \emph{vadhāre} again
\end{itemize}

\begin{tcolorbox}[breakable,colback=teal!4,colframe=teal!35!black,title={Before you move on}]
What does \gu{વધારે} mean? (\textquotedblleft{}More, too much\textquotedblright{}.) Say it once more.
\end{tcolorbox}
\section[\texorpdfstring{\gu{ટકો} (ṭako)}{ટકો (ṭako)}]{\gu{ટકો} (ṭako) — per cent}

\label{lesson:GU-C187-tako}



\begin{quote}
Before the new one: say the Gujarati for to go down, then the Gujarati for more.
\end{quote}

\subsection*{You'll want to know: \gu{ટકો}}
\textbf{\gu{ટકો}} — \emph{ṭako} — \textquotedblleft{}per cent\textquotedblright{}.

\begin{quote}
\textbf{\gu{પગાર} \gu{પાંચ} \gu{ટકા} \gu{વધે} \gu{છે}.} — \emph{pagār pā̃c ṭakā vadhe che} — \textquotedblleft{}The salary goes up by five per cent.\textquotedblright{}
\end{quote}

After a number it takes the plural \textbf{\gu{ટકા}}, \emph{ṭakā}.

\begin{cousinweb}
\textbf{\gu{ટકો}} (\emph{ṭako}) is a tadbhava. CDIAL lists it under Sanskrit \emph{ṭaṅka}, \textquotedblleft{}a weight; a stamped coin\textquotedblright{}, with the Gujarati sense \textquotedblleft{}two-paisā coin\textquotedblright{}; the sense \textquotedblleft{}per cent\textquotedblright{} is not in CDIAL (CDIAL 5426).
\end{cousinweb}

\begin{grammarlens}[title={What you've built}]
One more word for everyday life at A2.
\end{grammarlens}

\subsection*{Your turn}
\begin{itemize}
\raggedright
  \item \emph{Say it:} \emph{ṭako}
  \item \emph{Say it:} \emph{ṭako}, once more
  \item \emph{Say it:} \emph{pagār pā̃c ṭakā vadhe che} — The salary goes up by five per cent.
  \item \emph{Recall:} say the Gujarati for to go down, then the Gujarati for more, then say \emph{ṭako} again
\end{itemize}

\begin{tcolorbox}[breakable,colback=teal!4,colframe=teal!35!black,title={Before you move on}]
What does \gu{ટકો} mean? (\textquotedblleft{}Per cent\textquotedblright{}.) Say it once more.
\end{tcolorbox}
\section[\texorpdfstring{\gu{છૂટ} (chūṭ)}{છૂટ (chūṭ)}]{\gu{છૂટ} (chūṭ) — discount; freedom}

\label{lesson:GU-C187-chut}



\begin{quote}
Before the new one: say the Gujarati for more, then the Gujarati for per cent.
\end{quote}

\subsection*{You'll want to know: \gu{છૂટ}}
\textbf{\gu{છૂટ}} — \emph{chūṭ} — \textquotedblleft{}discount; freedom\textquotedblright{}.

\begin{quote}
\textbf{\gu{આ} \gu{દુકાનમાં} \gu{દસ} \gu{ટકા} \gu{છૂટ} \gu{છે}.} — \emph{ā dukānmā̃ das ṭakā chūṭ che} — \textquotedblleft{}This shop gives ten per cent off.\textquotedblright{}
\end{quote}

A feminine noun.

\begin{cousinweb}
\textbf{\gu{છૂટ}} (\emph{chūṭ}) is a tadbhava from the family that CDIAL lists under \emph{*kṣuṭyatē}, \textquotedblleft{}is released\textquotedblright{}, a reconstructed form; its Gujarati relatives there are \emph{chuṭvũ}, \textquotedblleft{}to get loose\textquotedblright{}, and \emph{chuṭũ}, \textquotedblleft{}loosened\textquotedblright{} (CDIAL 3707).
\end{cousinweb}

\begin{grammarlens}[title={What you've built}]
One more word for everyday life at A2.
\end{grammarlens}

\subsection*{Your turn}
\begin{itemize}
\raggedright
  \item \emph{Say it:} \emph{chūṭ}
  \item \emph{Say it:} \emph{chūṭ}, once more
  \item \emph{Say it:} \emph{ā dukānmā̃ das ṭakā chūṭ che} — This shop gives ten per cent off.
  \item \emph{Recall:} say the Gujarati for more, then the Gujarati for per cent, then say \emph{chūṭ} again
\end{itemize}

\begin{tcolorbox}[breakable,colback=teal!4,colframe=teal!35!black,title={Before you move on}]
What does \gu{છૂટ} mean? (\textquotedblleft{}Discount; freedom\textquotedblright{}.) Say it once more.
\end{tcolorbox}
